\section{Disparadores del Subsistema de Tendencias (Jesús)}


\subsection{Trigger: Procesar hashtags}

\textbf{Descripción:} En toda red social se lleva la cuenta de cuántas veces se ha mecionado un hashtag en todas las publicaciones de la misma. Sin embargo, este no es un procedimiento que se haga manualmente, sino que en cuanto que aparece un hashtag en una publicación, se actualiza la cuenta del mismo para ver qué temas son tendencia y cuales no. Por ello, es evidente la necesidad de implementarlo mediante un disparador.

\textbf{Funcionamiento:} Se ejecuta después de insertar (\texttt{AFTER INSERT}) en la tabla \texttt{PUBLICACION} (cuando algún usuario ha hecho una nueva publicación). Luego, se vé si en la descripción de dicha publicación aparece alguna cadena de texto que se corresponda con el formato de un hashtag. En ese caso, se comprueba si el hashtag ya existía (se actualiza su contador de menciones) o si no existía (se crea un nuevo hashtag con una mención). Finalmente se asocia dicha publicación con el/los hashtag/s correspondiente/s. 

\begin{lstlisting}[language=SQL, caption={Trigger TRG\_TEND\_PROCESAR\_HASHTAGS}]
CREATE OR REPLACE TRIGGER TRG_PROCESAR_HASHTAGS
AFTER INSERT ON PUBLICACION
FOR EACH ROW
DECLARE
    -- Variables para extraer los hashtags
    v_hashtag VARCHAR2(32);
    v_count   NUMBER;
    v_i       NUMBER := 1;
BEGIN
    -- 1. Buscamos todos los términos que empiecen por # en la descripción
    -- REGEXP_COUNT nos dice cuántos hashtags hay
    FOR i IN 1 .. REGEXP_COUNT(:NEW.DESCRIPCION, '#[[:alnum:]_]+') LOOP
        
        -- Extraemos el i-ésimo hashtag encontrado
        -- (Que empiecen por # y quedándonos con todos los caracteres alfanuméricos y los _)
        -- 1 -> buscar desde la primera posición
        -- i -> extraer el iésimo hashtag
        v_hashtag := REGEXP_SUBSTR(:NEW.DESCRIPCION, '#[[:alnum:]_]+', 1, i);

        -- Solo procesamos si el hashtag mide 32 o menos (según la especificación)
        IF LENGTH(v_hashtag) <= 32 THEN
            
            -- RS2.2: ¿Existe ya en la tabla HASHTAG?
            SELECT COUNT(*) INTO v_count FROM HASHTAG WHERE HASHTAG = v_hashtag;

            IF v_count > 0 THEN
                -- Si existe, actualizamos menciones
                UPDATE HASHTAG 
                SET MENCIONES = MENCIONES + 1 
                WHERE HASHTAG = v_hashtag;
            ELSE
                -- Si no existe, lo creamos con 1 mención
                INSERT INTO HASHTAG (HASHTAG, MENCIONES) 
                VALUES (v_hashtag, 1);
            END IF;

            -- Registramos la relación en la tabla CONTIENE_HASHTAG (Tabla 13)
            -- Usamos :NEW.IDPUBLICACION que es el ID que se acaba de generar
            INSERT INTO CONTIENE_HASHTAG (IDPUBLICACION, HASHTAG)
            VALUES (:NEW.IDPUBLICACION, v_hashtag);
            
        END IF;
    END LOOP;
END;
/
\end{lstlisting}


